\section{A two-wire boundary construction and its remaining estimate}
\label{sec:locc-boundary}

The failure of the ordered upper expression does not exclude other
interpolations between cut functionals. Bug\'ar and Vrana construct
upper functionals from standard operator geometric means\footnote{The
published Definition~2.1(G6), p.~2791, prints the convex-direction
inequality. For the standard scalar mean $G(a,b)=\sqrt{ab}$, the inputs
$(1,4)$ and $(4,1)$ at $t=1/2$ instead give $5/2>2$.
We use the usual concave geometric means, consistently with
Remark~2.2(iv) and the compression inequality in the proof of
Proposition~4.8, equation~(117), p.~2816. This is a local correction
of the printed direction.}
\cite[Theorem~4.5, Proposition~4.8 and Example~4.10]{BugarVrana2025}.
Here are the observables and moment used below. Put
\(\mathcal H=\bigotimes_{j=1}^k\mathcal H_j\) and
\(\mathcal H_S=\bigotimes_{j\in S}\mathcal H_j\).
For a cut \(b=S\mid S^c\), let \(P_\lambda^{\mathcal H_S}\) be
the Schur--Weyl isotypical projection indexed by the partition
\(\lambda\vdash n\) on \(\mathcal H_S^{\otimes n}\). If
\(H(\lambda/n)\) is the base-two Shannon entropy of its normalized
row lengths, define
\[
A_{b,n}=
\left.
\sum_{\lambda\vdash n}2^{nH(\lambda/n)}
P_\lambda^{\mathcal H_S}\otimes P_\lambda^{\mathcal H_{S^c}}
\right|_{\operatorname{Sym}^n(\mathcal H)}.
\]
These are the positive-definite observables of
\cite[equation~(118), p.~2817]{BugarVrana2025}.
Combining its defining equations~(1) and~(120), with
\(p=(1-\alpha)/\alpha\), gives
\[
g_\alpha(T)=
\lim_{n\to\infty}
\left\langle T^{\otimes n},
\mathbb G((A_{b,n}^p)_b)T^{\otimes n}
\right\rangle^{\alpha/n}.
\]
Thus the moment operator is
\(\mathbb G((A_{b,n}^p)_b)\).\footnote{Published equation~(121), p.~2817,
retains an additional outer \(1/p\) power. The powers cancel in the
defining equations~(1) and~(120); the proof of Proposition~5.1,
equation~(128), p.~2819, also uses the canceled expression.
This is a local displayed-formula correction.}
A choice of three paired cuts is exactly GHZ-normalized for every
\(0<\alpha<1\). Under the price-space theorem of
\cite[Proposition~C.3, p.~28]{Anthropic2026}, it suffices that the
spectral value of the two diagonal wires exceed the tree bound.
The resulting sufficient estimate allows a fixed positive deficit;
exact attainment remains a stronger special case that preserves
the prescribed normalized scale. We prove these criteria and the
normalization and limiting results behind them.
The required two-wire lower estimate remains open here.

All tensors in this section are over \(\mathbb C\), with a fixed number
\(k\geq2\) of parties and orthogonal local direct sums.
An abstract upper functional means the monotone, subadditive and
submultiplicative map of \cite[Definition~3.1]{BugarVrana2025}, with
\(g(0)=0\) and \(g(1)=1\). Normalization additionally means \(g(U_r)=r\)
for every \(r\). A normalized LOCC spectral point is additive,
multiplicative, and monotone under trace-nonincreasing LOCC.
Exact scalar order \(\alpha\) means
\(g(\sqrt p\,T)=p^\alpha g(T)\) for every \(p>0\).
We do not claim the stronger local-projection sum axiom of
\cite[Definition~4.1]{BugarVrana2025} for the growth-normalized upper
functional constructed below.
For a normalized upper functional $g$, write $\widetilde g$ for its
natural-number spectral regularization: it is the largest value at each
tensor attained by a normalized LOCC spectral point dominated by $g$.
Its precise formula is given in \eqref{eq:locc-regularization}.
The zero-order compactness lemma and the uniform-lower-bound criterion
are proved below before the detailed binary-mean calculations.

\subsection{The square candidate}

For convex cut weights \(\theta\), write \(g_\alpha=F^{\alpha,\mathbb G}\)
and \(d_\theta(x,y)=\sum_{b\text{ separating }x,y}\theta(b)\).

\begin{lemma}[Exact values on individual wires]
\label{lem:operator-mean-single-wire}
For every dimension-$q$ wire $X_{xy}^{(q)}$, with all unused parties
one-dimensional, and every $0<\alpha<1$,
\begin{equation}\label{eq:operator-mean-wire-value}
 g_\alpha(X_{xy}^{(q)})
 =q^{\alpha+(1-\alpha)d_\theta(x,y)}.
\end{equation}
\end{lemma}
\begin{proof}
Write $d=d_\theta(x,y)$ and $p=(1-\alpha)/\alpha$.
On the global symmetric $n$-copy space, every separating cut has the
same bipartite entropy observable $B_n$. Every nonseparating cut has a
one-dimensional side, so only the one-row partition occurs and its
observable is $I$. All the arguments of the mean therefore commute;
the scalar-weight property gives $G((A_n^{(b)})^p:b)=B_n^{pd}$.
Consequently
\[
 g_\alpha(X)=\lim_{n\to\infty}
    \langle X^{\otimes n},B_n^{pd}X^{\otimes n}\rangle^{\alpha/n}.
\]
For $d=0$, this is $\|X\|^{2\alpha}=q^\alpha$.
For $d>0$, put
$\beta=\alpha/(\alpha+(1-\alpha)d)$. Then
$(1-\beta)/\beta=pd$, so the same moment defines the bipartite
$\beta$-functional with exponent $\beta/n$.
Its value on the unnormalized $q$-wire is $q$
\cite[Lemma~4.6 and equation~(119)]{BugarVrana2025}.
Raising this limit to $\alpha/\beta$ gives the stated formula.
\end{proof}

\begin{lemma}[Paired cuts normalize GHZ tensors]\label{lem:paired-cut-ghz-normalization}
Suppose the four-party cut weights are supported on the three \(2|2\)
bipartitions.  For every standard operator geometric mean with these
weights and every \(0<\alpha<1\),
\[
 F^{\alpha,\mathbb G}(U_r)=r\qquad(r\geq1).
\]
\end{lemma}
\begin{proof}
Measure party four of \(U_r\) in the Fourier basis and correct the
outcome-dependent phase at party one.  This deterministic LOCC protocol
produces the three-party \(U_r\), with unchanged norm, and resets party
four to a one-dimensional space.  Every \(2|2\) cut then becomes a
nontrivial singleton cut on three parties.

The singleton Schur observables commute, so their mean is the weighted
product.  Their value on \(U_r\) is \(r\) for every positive \(\alpha\).
Here is a direct verification.  For fixed \(r\) and \(\delta>0\),
project each singleton onto Schur partitions whose entropy is at least
\(\log_2r-\delta\).  Its normalized \(n\)-copy marginal is \(r^{-n}I\).
The Schur--Weyl dimension bounds show that the probability discarded at
each singleton is at most a polynomial in \(n\) times \(2^{-n\delta}\).
The singleton projectors commute; hence their intersection has
probability \(1-o(1)\).  On that intersection the weighted product of
the observables raised to \(p=(1-\alpha)/\alpha\) is at least
\(2^{np(\log_2r-\delta)}\).
The defining moment, including \(\|U_r^{\otimes n}\|^2=r^n\), gives
\[
 F^{\alpha,\mathbb G}(U_r)
 \geq r^\alpha 2^{(1-\alpha)(\log_2r-\delta)}.
\]
Let \(\delta\downarrow0\).  Subadditivity gives the reverse bound \(r\).
LOCC monotonicity now bounds the four-party value below by the
three-party value \(r\), while subadditivity again bounds it above by \(r\).
\end{proof}

\begin{proposition}[Exact values on three-party wire products]
\label{prop:three-party-graph-locc}
Use any convex weights on the three four-party $2|2$ cuts and any
standard operator geometric mean with those weights. Let $G$ be a
product of wires, with edge dimensions $q_e$, supported on at most
three parties, and put
\[
 B_\alpha(G)=\prod_{e}q_e^{\alpha+(1-\alpha)d_\theta(e)}.
\]
For every $0<\alpha<1$ and integers $r,n\geq1$,
\begin{equation}\label{eq:three-party-graph-locc}
 g_\alpha(U_r\otimes G^{\otimes n})=r B_\alpha(G)^n,
 \qquad \widetilde g_\alpha(G)=B_\alpha(G).
\end{equation}
One dominated normalized LOCC spectral point attains all the
constituent wire ceilings for any fixed such $G$.
\end{proposition}
\begin{proof}
Choose a party outside the support of $G$. Fourier measurement of
its GHZ flag, followed by phase correction at another party, gives
the three-party $U_r\otimes G^{\otimes n}$ by deterministic LOCC
with unchanged norm. Every paired cut becomes a singleton cut, so
its Schur observables commute. The squared norm is
$r\prod_e q_e^n$, and the normalized marginal rank at cut $b$ is
$r\prod_{e\text{ crossing }b}q_e^n$. All these marginals are flat.
The simultaneous typical-Schur projection argument from the preceding
lemma, together with the matching bipartite-product upper bound
\cite[Proposition~5.1]{BugarVrana2025}, gives the exact output value
\[
 \left(r\prod_e q_e^n\right)^\alpha
 \prod_b\left(r\prod_{e\text{ crossing }b}q_e^n
         \right)^{(1-\alpha)\theta(b)}
 =r B_\alpha(G)^n.
\]
LOCC monotonicity gives the lower bound for the original tensor.
Paired-cut GHZ normalization, submultiplicativity, and the
individual-wire formula give the reverse bound.
The regularization formula~\eqref{eq:locc-regularization} now gives
$\widetilde g_\alpha(G)=B_\alpha(G)$.
A dominated maximizer at $G$ attains each constituent ceiling:
its product value equals the product of these ceilings, and each
individual logarithmic deficit is nonnegative.
\end{proof}

\begin{corollary}[A uniform baseline for the binary mean]
\label{cor:binary-uniform-baseline}
For the equal mean on $12|34$ and $14|23$, with
$W_q=\langle q\rangle_{13}\langle q\rangle_{24}$, one has
\begin{equation}\label{eq:binary-uniform-baseline}
 g_\alpha(U_r\otimes W_q)\geq r q^{1+\alpha},
 \qquad
 \widetilde g_\alpha(W_2)\geq2^{1+\alpha}
 \quad(r,q\geq1,\ 0<\alpha<1).
\end{equation}
\end{corollary}
\begin{proof}
Erasing the $24$ wire by Fourier measurement and phase correction
converts $W_q$ to $\sqrt q\,\langle q\rangle_{13}$ with unchanged
norm, while leaving the GHZ factor untouched. The preceding
proposition gives
$g_\alpha(U_r\otimes\langle q\rangle_{13})=rq$.
The scalar law gives the first inequality. Set $q=2^n$ and apply
\eqref{eq:locc-regularization} to obtain the second.
\end{proof}

The exponent surplus $\alpha$ in this baseline vanishes at the
zero-order boundary. The three-party equality does not extend here
to the four-party matching or to arbitrary three-party tensors.

\begin{corollary}[The baseline is optimal after removing a party]
\label{cor:inactive-party-output}
For the equal binary mean, if $S$ is an LOCC output of
$U_r\otimes W_q$ with one party one-dimensional, then
\[
 g_\alpha(S)\leq r q^{1+\alpha}.
\]
This bound is attained by such an output.
\end{corollary}
\begin{proof}
On $S$ the two paired cuts become two singleton cuts.
Proposition~5.1 of \cite{BugarVrana2025} bounds $g_\alpha(S)$ by
the geometric mean of their bipartite $\alpha$-functionals.
Each is LOCC-monotone. On the input, every singleton marginal is
flat of rank $rq$ and the squared norm is $rq^2$, so each such
functional has value
$(rq^2)^\alpha(rq)^{1-\alpha}=rq^{1+\alpha}$.
For equality, erase the $24$ wire and the fourth GHZ flag by the
Fourier protocols. The output is the three-party
$\sqrt q\,U_r\otimes\langle q\rangle_{13}$, whose value is
$q^\alpha rq$ by Proposition~\ref{prop:three-party-graph-locc}.
\end{proof}

This restriction on the preceding LOCC conversion is essential.
The corollary does not upper-bound the raw four-party functional.

\begin{proposition}[Exceeding the tree bound is sufficient]
\label{prop:locc-nontree-criterion}
Assume the price-space theorem of
\cite[Proposition~C.3]{Anthropic2026}.
Let \(H_i\) be normalized abstract upper functionals on the four-party
LOCC semiring, of positive exact scalar orders \(\eta_i\to0\).
Suppose, for some \(a\geq0\), that each cyclic side \(s\) satisfies
\[
 \limsup_i\log_2 H_i(\langle2\rangle_s)\leq a.
\]
Put \(W=\langle2\rangle_{13}\langle2\rangle_{24}\).
If
\[
 \liminf_i\log_2\widetilde H_i(W)>2a,
\]
then the unit square is a price.
\end{proposition}
\begin{proof}
Choose a dominated normalized LOCC spectral point \(h_i\leq H_i\)
attaining \(\widetilde H_i(W)\), using
\cite[Theorem~3.23]{BugarVrana2025}.
Reciprocal scalar tensors force its scalar order to be \(\eta_i\).
Lemma~\ref{lem:locc-zero-boundary} supplies a pointwise subnet limit
in the ordinary restriction spectrum.
Every side distance in its metric is at most \(a\).
Multiplicativity and the strict lower limit give
\[
 d_{13}+d_{24}>2a
 \geq\max\{d_{12}+d_{34},d_{14}+d_{23}\}.
\]
The limiting price is not a tree pseudometric. Apply
Theorem~\ref{thm:tree-price-dichotomy}.
\end{proof}

\begin{corollary}[Two equally weighted crossing cuts suffice]
\label{cor:binary-mean-square-sufficient}
Assume the complex price-space theorem of
\cite[Proposition C.3]{Anthropic2026}.
Use only the cuts \(12\mid34\) and \(14\mid23\), with weight
\(1/2\) each, in the operator geometric-mean construction, and
write its upper functional as \(g_\alpha\), \(0<\alpha<1\).
Put \(W=\langle2\rangle_{13}\langle2\rangle_{24}\).
If there are \(\alpha_i\to0\) and a fixed \(\delta>0\) such that
\begin{equation}\label{eq:binary-mean-positive-target}
g_{\alpha_i}(U_r\otimes W^{\otimes n})
\ge r\,2^{(1+\delta)n}
\qquad\text{for all integers }r,n\ge1,
\end{equation}
then the square is a price over \(\C\).
Equivalently, it suffices to replace the right side by
\(r4^n2^{-\varepsilon_i n}\) for nonnegative losses with
\(\limsup_i\varepsilon_i<1\), after discarding finitely many
indices.
\end{corollary}
\begin{proof}
Lemma~\ref{lem:paired-cut-ghz-normalization} gives
\(g_\alpha(U_r)=r\). A side is separated by exactly one of the
two cuts and a diagonal by both. Thus the binary side-wire
logarithms are \((1+\alpha)/2\), while each diagonal-wire
logarithm equals one. The product of the separate diagonal
values is consequently four.
By the regularization formula~\eqref{eq:locc-regularization},
\eqref{eq:binary-mean-positive-target} gives
\(\log_2\widetilde g_{\alpha_i}(W)\ge1+\delta\).
Proposition~\ref{prop:locc-nontree-criterion} applies with
\(a=1/2\).
For the loss formulation, choose
\(0<\delta<1-\limsup_i\varepsilon_i\) and omit an initial
finite segment; the converse uses the constant loss
\(\varepsilon_i=1-\delta\), reducing \(\delta\) first if
necessary. The two sufficient targets are therefore equivalent.
\end{proof}

The upper wire values in this two-cut construction do not form
a Euclidean square. The strict non-tree criterion does not
require them all to be attained. The unresolved estimate is
\eqref{eq:binary-mean-positive-target}, including its uniformity
in both the GHZ factor and the tensor power.

\begin{corollary}[A positive-deficit target for the paired-cut construction]
\label{cor:paired-cut-positive-deficit}
Assume the price-space theorem of
\cite[Proposition~C.3]{Anthropic2026}.
Label the square cyclically and put weight \(\sqrt2-1\) on each of
\(12|34\) and \(14|23\), and \(3-2\sqrt2\) on \(13|24\).
These convex cut weights have \(d_\theta=c e\), where
\(c=2-\sqrt2\).
For \(g_\alpha=F^{\alpha,\mathbb G}\), let
\[
 W=\langle2\rangle_{13}\langle2\rangle_{24},\qquad
 B_{\mathrm{diag},\alpha}
   =2^{\,2\alpha+2\sqrt2c(1-\alpha)}.
\]
If there are \(\alpha_i\to0\) and nonnegative \(\varepsilon_i\) with
\begin{equation}\label{eq:paired-cut-positive-deficit}
 \limsup_i\varepsilon_i<6\sqrt2-8
\end{equation}
such that
\[
 g_{\alpha_i}(U_r\otimes W^{\otimes n})
 \geq r B_{\mathrm{diag},\alpha_i}^{\,n}2^{-\varepsilon_i n}
 \qquad(r,n\geq1),
\]
then the square is a price.
The losses need not tend to zero; the endpoint in
\eqref{eq:paired-cut-positive-deficit} is not asserted.
\end{corollary}
\begin{proof}
A side is separated by cuts of total weight \(c\), and a diagonal
by cuts of total weight \(\sqrt2c=2\sqrt2-2\).
The exact single-wire formula gives the side logarithm
\(\alpha+c(1-\alpha)\), and \(B_{\mathrm{diag},\alpha}\)
is the product of the two diagonal wire values.
Lemma~\ref{lem:paired-cut-ghz-normalization} supplies raw normalization.
By Proposition~\ref{prop:locc-uniform-lower-bound}, the displayed
uniform estimate implies
\[
 \log_2\widetilde g_{\alpha_i}(W)
 \geq2\alpha_i+2\sqrt2c(1-\alpha_i)-\varepsilon_i.
\]
Since \(2c(\sqrt2-1)=6\sqrt2-8\), the lower limit exceeds \(2c\).
Apply Proposition~\ref{prop:locc-nontree-criterion} with \(a=c\).
No such uniform estimate is proved here.
\end{proof}

\begin{proposition}[Exact attainment preserves the normalized scale]\label{prop:diagonal-locc-criterion}
Assume the quadrilateral inequality for restriction-price metrics from
\cite[Proposition~C.3]{Anthropic2026}.
Let \(H_j\) be normalized abstract upper functionals on the four-party
LOCC semiring, of exact scalar orders \(\eta_j\to0\), whose binary-wire
values converge to \(2^{\lambda e_{xy}}\), with \(\lambda>0\) and \(e\)
the unit-side square metric.  Put
\(W=\langle2\rangle_{13}\langle2\rangle_{24}\).
If
\[
 \log_2\widetilde H_j(W)\longrightarrow2\sqrt2\,\lambda,
\]
then \(\lambda e\) is realized by one ordinary restriction spectral point.
\end{proposition}
\begin{proof}
Choose a dominated spectral point \(h_j\leq H_j\) attaining
\(\widetilde H_j(W)\), by \cite[Theorem~3.23]{BugarVrana2025}.
The exact scalar order is \(\eta_j\), by testing reciprocal scalar tensors.
The sum of the two diagonal wire logarithms tends to the sum of their
ceilings, so both nonnegative individual ceiling deficits tend to zero.
Lemma~\ref{lem:locc-zero-boundary} supplies a pointwise subnet limit
\(h_0\) in the restriction spectrum.  Its metric has both diagonals
\(\sqrt2\lambda\), and each side is at most \(\lambda\).
Apply the quadrilateral inequality only to this limiting restriction point:
\[
 4\lambda^2=d_{13}^2+d_{24}^2
 \leq d_{12}^2+d_{23}^2+d_{34}^2+d_{41}^2
 \leq4\lambda^2.
\]
Every side therefore equals \(\lambda\).
For a five-party centred square with the corresponding ten ceilings,
the same argument fixes the four corners; triangle equality on each
diameter then fixes both centre radii as well.
\end{proof}

\begin{corollary}[The exact-attainment special case]\label{cor:paired-cut-square-sufficient}
Use the family and notation of
Corollary~\ref{cor:paired-cut-positive-deficit}, with the same price-space
input.
If there are \(\alpha_i\downarrow0\) and \(\varepsilon_i\downarrow0\)
such that
\[
 g_{\alpha_i}(U_r\otimes W^{\otimes n})
 \geq r B_{\mathrm{diag},\alpha_i}^{\,n}2^{-\varepsilon_i n}
 \qquad(r,n\geq1),
\]
then the square metric is a price.
\end{corollary}
\begin{proof}
This is a special case of Corollary~\ref{cor:paired-cut-positive-deficit}.
For the stronger normalized conclusion, the raw normalized criterion of
Proposition~\ref{prop:locc-uniform-lower-bound}, applied to \(W\),
makes its spectral regularization approach the two-diagonal ceiling.
Proposition~\ref{prop:diagonal-locc-criterion} then applies with
\(\lambda=c\).
\end{proof}

\subsection{The zero-order limit and simultaneous values}

\begin{lemma}[The zero-order boundary]\label{lem:locc-zero-boundary}
Let $h_i$ be normalized LOCC spectral points with scalar orders
$0<\alpha_i\leq1$ tending to zero. There is a subnet converging pointwise
on all finite-dimensional $k$-party tensors to an ordinary normalized
restriction spectral point. Every pointwise cluster limit has this property.
\end{lemma}
\begin{proof}
Choose a set of representatives of finite Hilbert-space tensors. For each
nonzero tensor $T$, a finite rank decomposition supplies local maps $L_j$
with $T=(\bigotimes_jL_j)U_R$. Put
$c_j=\max\{1,\|L_j\|\}$ and $c=\prod_jc_j$. The maps $L_j/c_j$ are
contractions, so
\[
 h_i(T)\leq c^{2\alpha_i}R\leq c^2R.
\]
These coordinate bounds place the points in a product of compact intervals.
Tychonoff compactness gives a pointwise-convergent subnet, with limit $h_0$.
The additive and multiplicative identities and $h_0(U_r)=r$ pass to the
limit.

If $S=(\bigotimes_jA_j)T$, put
$d=\prod_j\max\{1,\|A_j\|\}$. Rescaling these maps to contractions gives
\[
 h_i(S)\leq d^{2\alpha_i}h_i(T).
\]
Hence $h_0(S)\leq h_0(T)$. This proves monotonicity under arbitrary
algebraic restrictions. A product-coordinate contraction extracts some
nonzero coefficient $z$ from each nonzero $T$, so
$h_i(T)\geq |z|^{2\alpha_i}$ and consequently $h_0(T)\geq1$.
Thus the limit is an ordinary normalized restriction spectral point.
An ordinary subsequence on the entire uncountable tensor semiring is
neither needed nor asserted.
\end{proof}

For a normalized upper functional $g$ on the LOCC semiring, let
$\widetilde g$ denote its $\mathbb N$-regularization in the sense of
\cite[Definition~3.13]{BugarVrana2025}. This is distinct from simply taking
roots of the values on tensor powers. The source proves
\begin{equation}\label{eq:locc-regularization}
 \widetilde g(T)
 =\lim_{n\to\infty}\inf_{r\geq1}
   \left(\frac{g(U_r\otimes T^{\otimes n})}{r}\right)^{1/n}
 =\max_{h\leq g}h(T),
\end{equation}
where the maximum is over normalized LOCC spectral points
\cite[Proposition~3.15, p.~2798, and Theorem~3.23, p.~2803]{BugarVrana2025}.
In particular this set of points is nonempty. The maximizer may depend on
$T$.

\begin{proposition}[A joint regularization deficit]
\label{prop:joint-locc-deficit}
For parameters $\alpha>0$ accumulating at zero, let $g_\alpha$ be LOCC
upper functionals of scalar order $\alpha$, normalized by
$g_\alpha(U_r)=r$ for every $r\geq1$. Fix nonzero tensors
$X_1,\ldots,X_\ell$, positive integers $m_1,\ldots,m_\ell$, and ceilings
\[
 g_\alpha(X_j)\leq a_{\alpha,j},\qquad
 a_{\alpha,j}\longrightarrow a_j>0.
\]
Set
\[
 Z=\bigotimes_{j=1}^{\ell}X_j^{\otimes m_j},\qquad
 B_\alpha=\prod_{j=1}^{\ell}a_{\alpha,j}^{m_j},\qquad
 \Delta_\alpha=\log_2B_\alpha-
                   \log_2\widetilde g_\alpha(Z).
\]
Then $\Delta_\alpha\geq0$. The condition
$\liminf_{\alpha\to0}\Delta_\alpha=0$ is equivalent to the existence of
$\alpha_i\to0$ and normalized LOCC spectral points
$h_i\leq g_{\alpha_i}$ such that $h_i(X_j)\to a_j$ for every $j$.
Either condition supplies a normalized restriction spectral point
$h_0$ with $h_0(X_j)=a_j$ simultaneously for all $j$.
\end{proposition}
\begin{proof}
Submultiplicativity gives $\widetilde g_\alpha(Z)\leq B_\alpha$.
By \eqref{eq:locc-regularization}, choose $h_\alpha\leq g_\alpha$ attaining
the value of $\widetilde g_\alpha$ on the single tensor $Z$. Then
\[
 \Delta_\alpha
 =\sum_jm_j\bigl(\log_2a_{\alpha,j}-\log_2h_\alpha(X_j)\bigr).
\]
Every summand is nonnegative. Thus
\[
 a_{\alpha,j}2^{-\Delta_\alpha/m_j}
 \leq h_\alpha(X_j)\leq a_{\alpha,j}.
\]
Along a sequence with $\Delta_\alpha\to0$, all the prescribed values
therefore converge together. Conversely any such choices satisfy
\[
 0\leq\Delta_{\alpha_i}
 \leq\sum_jm_j
       \bigl(\log_2a_{\alpha_i,j}-\log_2h_i(X_j)\bigr)
 \longrightarrow0.
\]

Every dominated point has scalar order exactly $\alpha$. Indeed domination
on $\sqrt p\,U_1$ and $p^{-1/2}U_1$, together with multiplicativity and
$h(U_1)=1$, gives
\[
 h(\sqrt p\,U_1)=p^\alpha.
\]
Multiplying by any tensor gives the required scalar law. The last conclusion
now follows from Lemma~\ref{lem:locc-zero-boundary}. Separate maximizers on
the $X_j$ would not establish this conclusion; the product $Z$ is essential.
\end{proof}

\begin{proposition}[The uniform normalized lower bound]
\label{prop:locc-uniform-lower-bound}
Under the hypotheses of Proposition~\ref{prop:joint-locc-deficit},
\(\Delta_\alpha\leq\varepsilon\) is equivalent to
\begin{equation}\label{eq:locc-uniform-lower-bound}
 g_\alpha(U_r\otimes Z^{\otimes n})
 \geq r B_\alpha^n2^{-\varepsilon n}\qquad(r,n\geq1).
\end{equation}
\end{proposition}
\begin{proof}
The sequence \(d_n=\inf_{r\geq1}g_\alpha(U_rZ^{\otimes n})/r\)
is positive, by any dominated normalized point, and submultiplicative,
using \(U_rU_s=U_{rs}\). Hence \eqref{eq:locc-regularization} equals
\[
 \widetilde g_\alpha(Z)=\inf_{n,r\geq1}
   \left(\frac{g_\alpha(U_rZ^{\otimes n})}{r}\right)^{1/n}.
\]
Comparing this infimum with \(B_\alpha2^{-\varepsilon}\) gives exactly
the stated condition.
\end{proof}


\subsection{A Schur complement for the equal two-cut mean}
\label{subsec:binary-mean}

For the equal-weight binary operator mean there is a useful elementary
block formula. We use the usual positive geometric mean
\[
 A\#B=A^{1/2}(A^{-1/2}BA^{-1/2})^{1/2}A^{1/2},
\]
as in \cite[Remark~2.2(ii), equation~(39)]{BugarVrana2025}.
The following is a block-matrix consequence of its classical defining
identities, rather than an additional asymptotic assertion.

\begin{lemma}[An involution and the geometric mean]
\label{lem:binary-mean-schur}
Let $A$ be a positive-definite operator on a finite-dimensional complex
Hilbert space, and let $R=R^*=R^{-1}$. Relative to the orthogonal
decomposition into the two eigenspaces of $R$, write
\[
 A=\begin{pmatrix}a&b\\b^*&c\end{pmatrix},
 \qquad S=a-bc^{-1}b^*.
\]
Then $A\#(RAR)$ commutes with $R$, and its block on the $+1$
eigenspace is
\begin{equation}\label{eq:binary-mean-schur}
 \bigl(A\#(RAR)\bigr)_+=a\#S.
\end{equation}
In particular, for $Rw=w$ its quadratic form on $w$ is
$\langle w,(a\#S)w\rangle$. If one eigenspace is zero-dimensional,
the corresponding one-block formula applies.
\end{lemma}
\begin{proof}
Put $B=RAR$ and $C=A\#B$. Symmetry and unitary covariance of the
geometric mean give $RCR=B\#A=C$, so $C=\operatorname{diag}(X,Y)$.
The geometric mean is the unique positive solution of
$CA^{-1}C=B$, as follows directly from its defining square root.
The block inverse formula gives $(A^{-1})_+=S^{-1}$. Taking the
$+1$ block in that equation yields
\[
 XS^{-1}X=a.
\]
The same characterization gives $X=S\#a=a\#S$.
No commutation of $a$ with $S$ is required.
\end{proof}

For the two-cut construction, put
$T=U_r\otimes\langle q\rangle_{13}\otimes\langle q\rangle_{24}$,
with each local Hilbert space identified with
$\mathbb C^r\otimes\mathbb C^q$.
On the global symmetric $n$-copy space, exchanging parties $2$ and $4$
is an orthogonal involution $R$ that fixes $T^{\otimes n}$ and
interchanges the Schur observables for $12|34$ and $14|23$.
For $p=(1-\alpha)/\alpha>0$, set
$A=(A_{12|34,n})^p$ and form its blocks as above.
Equation~\eqref{eq:binary-mean-schur} then rewrites the exact defining
moment as
\[
 \left\langle T^{\otimes n},
  \bigl((A_{12|34,n})^p\#(A_{14|23,n})^p\bigr)T^{\otimes n}\right\rangle
 =\left\langle T^{\otimes n},(a\#S)T^{\otimes n}\right\rangle.
\]
The Schur complement may be substantially smaller than $a$.
This formula supplies no lower estimate uniform in $r,n$ or
as $\alpha\to0$; the positive-deficit target remains a separate problem.

\subsection{A raw two-cut estimate from determinant polynomials}
\label{subsec:binary-raw-surplus}

The two-cut functional has a concrete lower bound on an individual
matching tensor. The estimate below concerns the raw upper functional;
it does not pass either of the two infima in its spectral regularization.
Its operator step is the common-vector lower bound of
\cite[Section~2.4, equation~(41)]{BugarVrana2025}; the determinant
polynomials provide a particular simultaneous Schur-support witness.

\begin{proposition}[A determinant-product witness]
\label{prop:binary-raw-surplus}
Let $g_\alpha$ be the Bug\'ar--Vrana upper functional obtained from
the equal-weight binary geometric mean of the observables for
$12|34$ and $14|23$, where $0<\alpha<1$. For $q\geq2$, put
\[
 W_q=\langle q\rangle_{13}\langle q\rangle_{24},\qquad d=q^2,
 \qquad
 h_d=H\left(\frac{d+1}{2d},\frac1{2d},\ldots,\frac1{2d}\right),
\]
where there are $d-1$ small coordinates and entropy has base two.
Then
\begin{equation}\label{eq:binary-raw-surplus}
 g_\alpha(W_q)\geq2^{(1-\alpha)h_d}.
\end{equation}
In particular,
\[
 \liminf_{\alpha\to0}\log_2g_\alpha(W_2)
 \geq h_4=3-\frac58\log_2 5>1.
\]
\end{proposition}
\begin{proof}
Put \(\mathcal H=(\mathbb C^q)^{\otimes4}\).
Let $D_0$ and $D_1$ be the determinant polynomials of the two
$d$ by $d$ adjacent flattenings. They have degree $d$ and take
values $\pm1$ on $W_q$. In the convention
$D_b(z)=\langle u_b,z^{\otimes d}\rangle$, their symmetric
coefficient vectors have norm one. Indeed, there are $d!$
distinct monomials and $d!$ distinct orderings of each, with
coefficients $\pm1/d!$ in the symmetric tensor.

For $m\geq1$, let $v_m$ represent $(D_0D_1)^m$ and put $N=2dm$.
Orthogonal symmetrization of $u_0^{\otimes m}\otimes u_1^{\otimes m}$
gives
\[
 \|v_m\|\leq1,\qquad
 |\langle v_m,W_q^{\otimes N}\rangle|=1.
\]
For either adjacent cut, multiplication by $D_b^m$ twists both
$\operatorname{GL}_d$ representations by the $m$th determinant
character. Thus each degree-$dm$ Schur partition $\mu$ in the
other factor is shifted to
$\lambda=\mu+(m,\ldots,m)$, with $d$ entries.
Every partition in $v_m$ therefore satisfies $\lambda_i\geq m$
and $|\lambda|=N$. The probability vector defining $h_d$
majorizes every probability vector whose $d$ coordinates are
at least $1/(2d)$. Hence all these Schur entropies are at least
$h_d$, simultaneously for both cuts.

Set $p=(1-\alpha)/\alpha$, and let $A,B$ be the $p$th powers
of the two Schur observables on $\operatorname{Sym}^N(\mathcal H)$.
They are at least $I$. If $P_m$ is the rank-one orthogonal
projection onto $v_m$, the common high-entropy support gives
\[
 A,B\ \geq\ I+(2^{Nph_d}-1)P_m.
\]
Monotonicity and idempotence of the geometric mean imply
$A\#B\geq2^{Nph_d}P_m$. Consequently
\[
 \langle W_q^{\otimes N},(A\#B)W_q^{\otimes N}\rangle
 \geq\frac{2^{Nph_d}}{\|v_m\|^2}
 \geq2^{Nph_d}.
\]
Taking the $\alpha/N$ power and then the defining limit along
$N=2dm$ proves \eqref{eq:binary-raw-surplus}.
\end{proof}

There are two separate limitations. First,
\[
 h_d=\log_2(2d)-\frac{d+1}{2d}\log_2(d+1),
 \qquad h_{4^n}=n+1+o(1).
\]
Since $W_2^{\otimes n}$ is locally equivalent to $W_{2^n}$,
the estimate supplies no fixed positive excess over the tree
threshold per binary matching copy. Power-multiplicativity of
$g_\alpha$ has not been proved. Second, it does not supply the
lower bound uniform in the independent GHZ factor $U_r$.
Thus the strict lower bound for the raw value at $W_2$ alone
does not establish a square price.

\begin{corollary}[A necessary loss under spectral regularization]
\label{cor:binary-regularization-loss}
Assume the price-space theorem of
\cite[Proposition~C.3]{Anthropic2026}, and let
$\widetilde g_\alpha$ be the $\mathbb N$-regularization of the
equal binary-cut functional above. Then
\[
 \limsup_{\alpha\to0}\log_2\widetilde g_\alpha(W_2)\leq\sqrt2,
 \qquad
 \liminf_{\alpha\to0}
 \left(\log_2g_\alpha(W_2)-\log_2\widetilde g_\alpha(W_2)\right)
 \geq h_4-\sqrt2>0.
\]
If the square is not a price, the first upper bound improves to $1$
and the second lower bound improves to $h_4-1$.
\end{corollary}
\begin{proof}
Choose a sequence $\alpha_i\to0$ approaching the upper limit in
the first assertion. Theorem~3.23 of \cite{BugarVrana2025} supplies
dominated normalized LOCC spectral points $f_i\leq g_{\alpha_i}$
with $f_i(W_2)=\widetilde g_{\alpha_i}(W_2)$.
Their scalar orders are $\alpha_i$, by testing reciprocal scalar
tensors. By Lemma~\ref{lem:locc-zero-boundary}, a pointwise subnet
converges to an ordinary normalized restriction spectral point $f$.
The individual-wire formula~\eqref{eq:operator-mean-wire-value}
implies that each cyclic side in its price metric has length at most
$1/2$. Multiplicativity gives the limiting objective $d_{13}+d_{24}$.
The CAT(0) quadrilateral inequality gives
\[
 (d_{13}+d_{24})^2
 \leq2(d_{13}^2+d_{24}^2)
 \leq2(d_{12}^2+d_{23}^2+d_{34}^2+d_{41}^2)
 \leq2.
\]
If the square is not a price, Theorem~\ref{thm:tree-price-dichotomy}
makes this a tree pseudometric, whose four-point condition instead
gives
$d_{13}+d_{24}\leq\max\{d_{12}+d_{34},d_{14}+d_{23}\}\leq1$.
Combine these upper limits with Proposition~\ref{prop:binary-raw-surplus}.
Finally, $5^5<2^{12}$ implies $h_4>3/2>\sqrt2$.
\end{proof}

This loss is forced even without deciding whether the square is a
price. The conclusion does not determine whether the loss occurs
in the tensor-power infimum, in the GHZ-factor infimum, or in both.


We use the \emph{asymptotic subrank}
$\widetilde Q(T)=\lim_{n\to\infty}Q(T^{\otimes n})^{1/n}$,
where $Q(T)$ is the largest $q$ for which the unit tensor $U_q$
is a restriction of $T$ on the same party set. Its value is bounded
above by every ordinary normalized spectral value of $T$.
The required four-cycle calculation is proved in
Lemma~\ref{lem:cycle4-subrank}.

\begin{corollary}[The graph-subrank face and its exact baseline]
\label{cor:binary-graph-face}
Every ordinary normalized four-party spectral point whose four cyclic
side distances are at most $1/2$ has all four sides equal to $1/2$,
and satisfies
\[
 d_{13}+d_{24}\geq1.
\]
Equality holds precisely when its finite wire metric is the star with
pendant lengths $(a,b,a,b)$ for $a,b\geq0$, $a+b=1/2$.
For the equal binary-cut mean,
\begin{equation}\label{eq:binary-graph-boundary}
 \liminf_{\alpha\to0}\log_2\widetilde g_\alpha(W_2)\geq1.
\end{equation}
Assuming the price-space theorem, if the square is not a price,
this is an actual limit equal to one.
\end{corollary}
\begin{proof}
Lemma~\ref{lem:cycle4-subrank} gives asymptotic subrank $4$ for the
binary four-cycle. Its spectral value is two raised to the sum of
its four edge distances. Applying this to cycle $1234$ forces all
four cyclic sides to equal $1/2$. Apply the same lemma to the
relabeled cycle $1243$. Its edge sum is
\[
 d_{12}+d_{24}+d_{43}+d_{31}=1+d_{13}+d_{24}\geq2,
\]
which proves the diagonal lower bound.
If equality holds, set $a=d_{13}/2$ and $b=d_{24}/2$.
The star with pendants $(a,b,a,b)$ has exactly the given six distances.
Conversely such a star has diagonal sum one.

Choose dominated maximizers for the binary mean along a sequence
approaching the lower limiting objective. Every zero-order compactness
limit has side ceilings $1/2$, so the preceding argument bounds its
diagonal sum below by one. Multiplicativity identifies that sum with
the corresponding limit of $\log_2\widetilde g_\alpha(W_2)$.
This proves \eqref{eq:binary-graph-boundary}.
If the square is not a price, Theorem~\ref{thm:tree-price-dichotomy}
makes each limiting metric a tree. Its four-point rule bounds the
same diagonal sum above by one. The bounded objectives therefore
have only one accumulation value, giving the complete limit.
\end{proof}

This argument uses only the two four-cycle subrank inequalities;
no geodesic extension is needed for the baseline or its equality case.
It does not classify the spectral point on other tensors or provide
any strict surplus above one.

\subsection{General normalization by unit growth}

\begin{proposition}[Normalization by unit growth]\label{prop:ghz-growth-normalization}
Let \(g\) be a finite LOCC-monotone upper functional, positive on nonzero tensors, with \(g(1)=1\), \(g(0)=0\), and exact scalar order \(\alpha\in(0,1]\).  Define
\[
 \beta=\lim_{r\longrightarrow\infty}
       \frac{\log g(U_r)}{\log r}.
\]
The limit exists, \(\alpha\leq\beta\leq1\), and
\begin{equation}\label{eq:ghz-growth-normalization}
 H(T)=\left(\lim_{m\longrightarrow\infty}g(T^{\otimes m})^{1/m}\right)^{1/\beta}.
\end{equation}
is a normalized upper functional on the LOCC preordered semiring, of exact scalar order \(\alpha/\beta\).  In particular, \(H(U_r)=r\).
For the operator-mean family with cut weights \(\theta\), put
\(a_{xy}=\alpha+(1-\alpha)d_\theta(x,y)\).  Then
\[
 H(\langle q\rangle_{xy})=q^{a_{xy}/\beta},
 \qquad \beta\geq\max_{x\ne y}a_{xy}.
\]
\end{proposition}
\begin{proof}
Write \(a_r=g(U_r)\).  Restriction by local contractions makes \(a_r\) nondecreasing.  Submultiplicativity gives \(a_{rs}\leq a_ra_s\).  Subadditivity gives \(a_r\leq r\), whereas deterministic LOCC erasure \(U_r\to\sqrt r\,1\) gives \(a_r\geq r^\alpha\).  Set
\(\beta_*=\inf_{r\geq2}\log a_r/\log r\).
For fixed \(r\geq2\) and \(r^k\leq N<r^{k+1}\),
\(a_N\leq a_r^{k+1}\).  Hence the upper limit of
\(\log a_N/\log N\) is at most \(\log a_r/\log r\).  Taking the infimum proves existence of the limit and equality with \(\beta_*\).

Submultiplicativity and Fekete's lemma give
\(F(T)=\lim_m g(T^{\otimes m})^{1/m}\).
Extracting a nonzero product coefficient \(z\) by local contractions shows
\(g(T^{\otimes m})\geq |z|^{2\alpha m}\); thus \(F(T)>0\) for nonzero \(T\).
Monotonicity and submultiplicativity pass to \(F\).  Its scalar order is \(\alpha\), and
\(F(U_r)=r^\beta\).

It remains to check subadditivity after raising \(F\) to \(1/\beta\).
Fix \(p>\beta\), \(a>F(S)\), and \(b>F(T)\).  There are finite constants
\(C,C_S,C_T\) such that, for all integers \(\ell\geq1\) and \(j\geq0\),
\[
 g(U_\ell)\leq C\ell^p,\qquad
 g(S^{\otimes j})\leq C_Sa^j,\qquad
 g(T^{\otimes j})\leq C_Tb^j.
\]
The direct-sum binomial expansion therefore gives
\[
 g((S\oplus T)^{\otimes m})
 \leq CC_SC_T\sum_{j=0}^m\binom mj^p a^jb^{m-j}
 \leq CC_SC_T(m+1)(a^{1/p}+b^{1/p})^{pm}.
\]
The last inequality bounds each summand using the ordinary binomial theorem.
Take \(m\)-th roots and limits, then let \(a\downarrow F(S)\),
\(b\downarrow F(T)\), and \(p\downarrow\beta\).  This proves
\(F(S\oplus T)^{1/\beta}\leq F(S)^{1/\beta}+F(T)^{1/\beta}\).
The zero cases are immediate.  The other upper-functional axioms and the scalar order now follow.

The exact wire evaluation from Lemma~\ref{lem:operator-mean-single-wire}
holds for every grouped tensor power, since
\(\langle q\rangle_{xy}^{\otimes m}\cong\langle q^m\rangle_{xy}\).
This proves the displayed value of \(H\).
Finally \(U_q\to\langle q\rangle_{xy}\) is a deterministic LOCC conversion,
so \(g(U_q)\geq q^{a_{xy}}\); its growth limit gives the bound on \(\beta\).
\end{proof}

\begin{proposition}[The raw lower bound after normalization]\label{prop:raw-locc-deficit}
Use Proposition~\ref{prop:ghz-growth-normalization}.  Let \(X_1,\ldots,X_s\)
be wires, let \(Z=\bigotimes_jX_j\), and let \(B=\prod_jg(X_j)\).
Let \(\widetilde H\) be the spectral regularization of the normalized upper functional \(H\), and put
\[
 \Delta=\log_2 B-\beta\log_2\widetilde H(Z)\geq0.
\]
For every \(\varepsilon\geq0\), the inequality \(\Delta\leq\varepsilon\)
is equivalent to
\begin{equation}\label{eq:raw-locc-uniform}
 g(U_r\otimes Z^{\otimes n})
 \geq r^\beta B^n2^{-\varepsilon n}
 \qquad\text{for all integers }r,n\geq1.
\end{equation}
Its \(n=0\) version is automatic.
For the operator-mean family with fixed nonzero cut metric, suppose in
addition that the list contains a binary wire for every unordered pair.
If \(\alpha_i\to0\) and \(\Delta_i\to0\), a subnet of jointly maximizing
dominated spectral points converges to an ordinary restriction spectral
point with wire metric \(d_\theta/\beta_0\), where a subsequence has
\(\beta_i\to\beta_0\in[d_{\max},1]\).
\end{proposition}
\begin{proof}
By \cite[Proposition~3.15 and Theorem~3.23]{BugarVrana2025},
\[
 \widetilde H(Z)
 =\inf_{n\geq1}\inf_{r\geq1}
       \left(\frac{H(U_r\otimes Z^{\otimes n})}{r}\right)^{1/n}.
\]
The limit in the source formula equals this infimum because the
GHZ-assisted sequence is submultiplicative.
Also \(H(W)\leq g(W)^{1/\beta}\).
Consequently \(\Delta\leq\varepsilon\) implies
\[
 g(U_rZ^n)\geq H(U_rZ^n)^\beta
 \geq r^\beta\widetilde H(Z)^{\beta n}
 \geq r^\beta B^n2^{-\varepsilon n}.
\]
Conversely apply \eqref{eq:raw-locc-uniform} with \(r^m,nm\) and take the
\(1/(\beta m)\)-th power and the limit.  Since
\((U_rZ^n)^{\otimes m}\cong U_{r^m}Z^{nm}\), this gives
\(H(U_rZ^n)\geq rB^{n/\beta}2^{-\varepsilon n/\beta}\).
The regularization formula gives \(\Delta\leq\varepsilon\).
For \(n=0\), the assertion follows from
\(g(U_r)\geq r^\beta\), the infimum characterization of \(\beta\).

For this last assertion, the exact normalized wire ceilings are
\(g(X_j)^{1/\beta}\).
The one-product maximization argument of
Proposition~\ref{prop:joint-locc-deficit} supplies dominated points
simultaneously approaching them when \(\Delta\to0\).
Since \(\beta\geq\alpha+(1-\alpha)d_{\max}\) and \(d_{\max}>0\),
the scalar orders \(\alpha/\beta\) tend to zero.
Lemma~\ref{lem:locc-zero-boundary} then applies, and the wire exponents
converge to \(d_\theta/\beta_0\).
\end{proof}

\begin{lemma}[Unit factors after power regularization]\label{lem:power-upper-unit-regularization}
Let \(H\) be a normalized abstract upper functional on a semiring with a
Strassen preorder. Suppose \(H(y^n)=H(y)^n\) for every semiring element
\(y\) and every \(n\geq1\). Then, for every \(x\),
\[
 \widetilde H(x)=\inf_{r\geq1}\frac{H(rx)}r
 =\lim_{r\longrightarrow\infty}\frac{H(rx)}r.
\]
Consequently the raw bound in \eqref{eq:raw-locc-uniform} is equivalently
tested only on
\[
 g(U_{r^n}Z^{\otimes n})
 \geq r^{\beta n}B^n2^{-\varepsilon n}
 \qquad(r,n\geq1).
\]
The tolerance must remain uniform in \(r\).
\end{lemma}
\begin{proof}
For each \(n\), subadditivity gives
\(L_n=\lim_{r\to\infty}H(rx^n)/r=\inf_rH(rx^n)/r\).
Along \(r=s^n\), power multiplicativity gives
\(H(s^nx^n)/s^n=(H(sx)/s)^n\), so \(L_n=L_1^n\).
The formula in \cite[Proposition~3.15]{BugarVrana2025} gives
\(\widetilde H(x)=L_1\).
The growth-normalized \(H\) is power-multiplicative by its defining limit.
For fixed \(r\), Fekete's lemma gives
\[
 H(U_rZ)^\beta=\inf_{n\geq1}
             g(U_{r^n}Z^{\otimes n})^{1/n}.
\]
Insert this in the displayed unit-factor formula.
This argument does not assume that the raw \(g\) is power-multiplicative.
\end{proof}


For a general operator-mean square ray \(d_\theta=c e\), apply the
growth-normalized \(H_\alpha\) to \(W=\langle2\rangle_{13}\langle2\rangle_{24}\).
Its two-diagonal deficit is
\[
 \Delta_{\mathrm{diag},\alpha}
 =2\alpha+2\sqrt2c(1-\alpha)
       -\beta_\alpha\log_2\widetilde H_\alpha(W).
\]
The raw criterion of Proposition~\ref{prop:raw-locc-deficit}, with
\(Z=W\) and \(B_\alpha=2^{2\alpha+2\sqrt2c(1-\alpha)}\), is exact.
If this deficit tends to zero with \(\alpha\), a subsequence
\(\beta_\alpha\to\beta_0>0\) gives scalar order
\(\alpha/\beta_\alpha\to0\) and scale \(c/\beta_0\) in
Proposition~\ref{prop:diagonal-locc-criterion}.

The established flat-marginal evaluation in
\cite[Corollary~5.5, p.~2821]{BugarVrana2025} supplies equality for
GHZ tensors, graph tensors and their products only for
\(\alpha\in[1/2,1)\). Its extension theorem gives global dominated
LOCC points preserving that generated subsemiring
\cite[Corollary~3.24, p.~2805]{BugarVrana2025}.
This parameter range does not approach zero.
Our normalization results remove a separate unit-normalization
prerequisite; they do not establish the required positive-\(n\)
two-wire estimate. Individual-wire values or finite-copy raw lower
bounds do not, by themselves, control the spectral regularization.
